% Modular TeX snippet for the 253b final paper.
% Mathematical prerequisites: algebraic topology, symplectic geometry, and Lie groups for TQFT

\providecommand{\dd}{\mathrm{d}}
\providecommand{\Tr}{\operatorname{Tr}}
\providecommand{\tr}{\operatorname{tr}}
\providecommand{\CS}{\operatorname{CS}}
\providecommand{\Hom}{\operatorname{Hom}}
\providecommand{\U}{\operatorname{U}}
\providecommand{\SU}{\operatorname{SU}}
\providecommand{\SO}{\operatorname{SO}}
\providecommand{\Spin}{\operatorname{Spin}}
\providecommand{\Lk}{\operatorname{Lk}}
\providecommand{\PD}{\operatorname{PD}}
\providecommand{\Hol}{\operatorname{Hol}}
\providecommand{\Bord}{\operatorname{Bord}}
\providecommand{\Vect}{\operatorname{Vect}}

\section{Topology and symplectic geometry for Chern--Simons theory}
\label{sec:topology-review}

Low-dimensional quantum field theory depends on global field structure as much as local differential equations. Flat gauge fields, though locally pure gauge, carry nontrivial holonomy. Bundle topology distinguishes $\SU(2)$ from $\SO(3)$ gauge theory. The moduli space of flat connections is symplectic. Functorial TQFT encodes gluing in bordism categories. We collect the necessary mathematical structures.

\subsection{Homotopy and fundamental groups}

\subsubsection{Homotopy equivalence}

Let $X$ and $Y$ be topological spaces. Maps $f : X \to Y$ and $g : Y \to X$ form a \emph{homotopy equivalence} if
\begin{equation}
  g \circ f \simeq \operatorname{id}_X,
  \qquad
  f \circ g \simeq \operatorname{id}_Y,
\end{equation}
where $\simeq$ denotes homotopy. A \emph{homotopy} between maps $f_0,f_1 : X \to Y$ is a continuous map $H : X \times [0,1] \to Y$ with $H(x,0)=f_0(x)$ and $H(x,1)=f_1(x)$. Homotopy equivalence is the correct notion of sameness for topological field theory because metric deformations do not alter the partition function.

\subsubsection{The fundamental group}

A \emph{path} from $x_0$ to $x_1$ in $X$ is a continuous map $\gamma : [0,1] \to X$ with $\gamma(0)=x_0$ and $\gamma(1)=x_1$. A \emph{loop} based at $x_0$ is a path with $x_1=x_0$. Two based loops are homotopic rel endpoints if there is a homotopy through loops keeping endpoints fixed. The \emph{fundamental group}
\begin{equation}
  \pi_1(X,x_0)
\end{equation}
is the set of based homotopy classes of loops with group operation given by concatenation $[\gamma_1] \cdot [\gamma_2] = [\gamma_1 * \gamma_2]$. The constant loop is the identity; the reversed loop $\overline\gamma(t)=\gamma(1-t)$ gives the inverse.

For a flat connection, parallel transport around loops depends only on homotopy class. Hence flat gauge fields define representations
\begin{equation}
  \rho : \pi_1(M) \to G.
\end{equation}
This is the basic bridge between topology and gauge theory.

\subsubsection{\texorpdfstring{Computation of $\pi_1(S^1)$ and $\pi_1(T^2)$}{Computation of pi1(S1) and pi1(T2)}}

The universal covering map
\begin{equation}
  p : \mathbb{R} \to S^1,
  \qquad
  p(t)=e^{2\pi i t}
\end{equation}
lifts every loop $\gamma : [0,1] \to S^1$ with $\gamma(0)=\gamma(1)=1$ to a path $\widetilde\gamma : [0,1] \to \mathbb{R}$ with $\widetilde\gamma(0)=0$. Since $p(\widetilde\gamma(1))=1$, the endpoint $\widetilde\gamma(1)$ is an integer. The \emph{degree}
\begin{equation}
  \deg(\gamma)=\widetilde\gamma(1) \in \mathbb{Z}
\end{equation}
is the winding number. Homotopic loops lift to homotopic paths with the same endpoint; concatenation adds endpoints; every integer occurs. Hence
\begin{equation}
  \pi_1(S^1) \cong \mathbb{Z}.
\end{equation}

For the torus $T^2 = S^1 \times S^1$,
\begin{equation}
  \pi_1(T^2) \cong \pi_1(S^1) \times \pi_1(S^1) \cong \mathbb{Z}^2.
\end{equation}
A flat abelian connection on the torus is controlled by two holonomies around meridian and longitude.

\subsubsection{Higher homotopy groups}

For $n \ge 1$, define
\begin{equation}
  \pi_n(X,x_0) = [ (S^n,s_0),(X,x_0) ],
\end{equation}
the set of based homotopy classes of maps from $S^n$ to $X$. For $n>1$, $\pi_n(X,x_0)$ is abelian. In Chern--Simons theory the integer
\begin{equation}
  \int_M \tr(g^{-1}\dd g)^{\wedge 3},
  \qquad
  g : M^3 \to G
\end{equation}
is controlled by the homotopy class $[g]\in\pi_3(G)$. For compact simple simply connected $G$,
\begin{equation}
  \pi_3(G) \cong \mathbb{Z},
\end{equation}
which is why the Chern--Simons level is quantized.

\subsubsection{Long exact sequence of a fibration}

If $F \hookrightarrow E \xrightarrow{p} B$ is a fibration, there is a long exact sequence
\begin{equation}
  \cdots \to \pi_n(F) \to \pi_n(E) \to \pi_n(B)
  \xrightarrow{\partial}
  \pi_{n-1}(F) \to \cdots \to \pi_0(B).
\end{equation}

The fibration
\begin{equation}
  \SU(n-1) \hookrightarrow \SU(n) \to S^{2n-1}
\end{equation}
sends a matrix to its first column. The stabilizer of a unit vector is $\SU(n-1)$. Since $\pi_3(S^{2n-1})=0$ for $n\ge 3$, exactness gives
\begin{equation}
  \pi_3(\SU(n-1)) \xrightarrow{\sim} \pi_3(\SU(n))
\end{equation}
for $n\ge 3$. Since $\pi_3(\SU(2))=\mathbb{Z}$, induction yields
\begin{equation}
  \pi_3(\SU(n)) \cong \mathbb{Z}
  \qquad
  (n\ge 2).
\end{equation}
This is the topological statement underlying the winding-number formula in nonabelian Chern--Simons theory.

\subsection{Bundles, connections, and characteristic classes}

\subsubsection{Principal bundles and transition functions}

A \emph{principal $G$-bundle} is a map $\pi : P \to B$ with free right $G$-action, simply transitive on fibers, and locally trivial. On an open cover $\{U_i\}$ of $B$, the bundle is described by transition functions
\begin{equation}
  g_{ij} : U_i \cap U_j \to G
\end{equation}
satisfying the cocycle condition
\begin{equation}
  g_{ij} g_{jk} g_{ki} = e
\end{equation}
on triple overlaps. Bundle isomorphism classes are represented by Čech $1$-cocycles with values in $G$.

For a topological group $G$, principal $G$-bundles over $X$ are classified by homotopy classes of maps
\begin{equation}
  [X,BG],
\end{equation}
where $BG$ is the classifying space. There is a universal bundle $EG \to BG$, and every principal $G$-bundle over $X$ is pulled back from this by a classifying map $f : X \to BG$.

\subsubsection{Connections and curvature}

A \emph{connection} on $P \to M$ is a $\mathfrak g$-valued $1$-form $A \in \Omega^1(P;\mathfrak g)$ reproducing generators on vertical vectors and equivariant under the $G$-action. Locally $A$ is a gauge potential. The \emph{curvature} is
\begin{equation}
  F = \dd A + A \wedge A.
\end{equation}
Under gauge transformation $g : M \to G$,
\begin{equation}
  A \mapsto g^{-1}Ag + g^{-1}\dd g,
  \qquad
  F \mapsto g^{-1}Fg.
\end{equation}
The Bianchi identity is
\begin{equation}
  D_A F = \dd F + [A,F]=0.
\end{equation}

\subsubsection{Chern--Weil theory and transgression}

Let $P$ be an invariant polynomial on $\mathfrak g$. The differential form $P(F)$ is closed, and its de Rham cohomology class is independent of connection. Hence it defines a characteristic class of the bundle. For complex vector bundles the Chern character is
\begin{equation}
  \operatorname{ch}(E)=\tr\exp\!\left(\frac{iF}{2\pi}\right).
\end{equation}

For a connection $A$, define
\begin{equation}
  \omega_{\CS}(A)=\tr\!\left(A\wedge \dd A + \frac{2}{3}A\wedge A\wedge A\right).
\end{equation}
Direct calculation gives
\begin{equation}
  \dd\omega_{\CS}(A)=\tr(F\wedge F).
\end{equation}
Thus the three-dimensional Chern--Simons form is a transgression form for the four-dimensional characteristic class. This relation is the local algebraic origin of anomaly descent.

\subsection{Homology, cohomology, and linking}

\subsubsection{Singular homology and de Rham cohomology}

The singular chain group $C_n(X;\mathbb{Z})$ is the free abelian group on continuous maps $\sigma : \Delta^n \to X$. The boundary operator $\partial : C_n(X) \to C_{n-1}(X)$ satisfies $\partial^2=0$. The homology groups are
\begin{equation}
  H_n(X;\mathbb{Z}) = \ker \partial / \operatorname{im} \partial.
\end{equation}
Dually, cochains are $C^n(X;\mathbb{Z}) = \Hom(C_n(X),\mathbb{Z})$ with coboundary $\delta$, and
\begin{equation}
  H^n(X;\mathbb{Z}) = \ker \delta / \operatorname{im} \delta.
\end{equation}

For a smooth manifold $M$, let $\Omega^k(M)$ be the smooth $k$-forms with exterior derivative $\dd : \Omega^k(M) \to \Omega^{k+1}(M)$. The de Rham cohomology groups are
\begin{equation}
  H^k_{\mathrm{dR}}(M) = \ker(\dd : \Omega^k \to \Omega^{k+1}) / \operatorname{im}(\dd : \Omega^{k-1} \to \Omega^k).
\end{equation}
The de Rham theorem states
\begin{equation}
  H^k_{\mathrm{dR}}(M) \cong H^k(M;\mathbb{R}).
\end{equation}
This is the bridge on which Chern--Weil theory is built.

\subsubsection{Stokes' theorem and transgression}

Stokes' theorem says that for a compact oriented manifold with boundary,
\begin{equation}
  \int_M \dd\omega = \int_{\partial M} \omega.
\end{equation}
Whenever a topological density is written as an exact form $P(F)=\dd Q(A)$, Stokes' theorem turns bulk information into boundary information. This is the mathematical core of anomaly inflow.

\subsubsection{Poincaré duality and linking}

For a closed oriented $n$-manifold $M$, Poincaré duality gives
\begin{equation}
  H^k(M;R) \cong H_{n-k}(M;R).
\end{equation}
A codimension-$k$ oriented submanifold $N \subset M$ determines a cohomology class
\begin{equation}
  \PD[N] \in H^k(M;R)
\end{equation}
characterized by
\begin{equation}
  \int_N \iota^*\eta = \int_M \eta \wedge \PD[N]
\end{equation}
for all closed $(n-k)$-forms $\eta$. Worldlines are naturally homology classes; gauge fields are naturally cohomology classes. Poincaré duality converts between them.

Linking numbers in dimension three can be expressed cohomologically. If $C_1,C_2 \subset M^3$ are disjoint oriented closed curves and $S$ is a Seifert surface with $\partial S=C_1$, then
\begin{equation}
  \Lk(C_1,C_2)= [S]\cdot [C_2].
\end{equation}
Equivalently, if $\alpha$ is a one-form with $\dd\alpha = \PD[C_1]$, then
\begin{equation}
  \Lk(C_1,C_2)=\int_{C_2} \alpha.
\end{equation}
This is the topological content of abelian Wilson-loop linking formulas in Chern--Simons theory.

\subsection{Lie groups and homotopy}

A \emph{Lie group} is a smooth manifold and a group with smooth multiplication and inversion. From the topological viewpoint, a Lie group is a highly structured space with computable homotopy type. Both viewpoints are indispensable in TQFT.

\subsubsection{\texorpdfstring{$\SU(2)$ is the three-sphere}{SU(2) is the three-sphere}}

Every element of $\SU(2)$ has the form
\begin{equation}
  \begin{pmatrix}
  a & b \\
  -\overline b & \overline a
  \end{pmatrix}
\end{equation}
with $|a|^2+|b|^2=1$. Writing $a=x_1+ix_2$ and $b=x_3+ix_4$, this becomes $x_1^2+x_2^2+x_3^2+x_4^2=1$. Hence
\begin{equation}
  \SU(2) \cong S^3.
\end{equation}
Therefore
\begin{equation}
  \pi_1(\SU(2))=0,
  \qquad
  \pi_3(\SU(2))\cong \pi_3(S^3)\cong \mathbb{Z}.
\end{equation}
This is the prototype of all winding-number statements.

\subsubsection{\texorpdfstring{$\SO(3)$ and its double cover}{SO(3) and its double cover}}

The surjective homomorphism $\SU(2) \to \SO(3)$ with kernel $\{\pm I\}$ gives
\begin{equation}
  \SO(3) \cong \SU(2)/\{\pm I\} \cong S^3/\{\pm 1\} \cong \mathbb{R}P^3.
\end{equation}
Since $S^3 \to \mathbb{R}P^3$ is a double cover,
\begin{equation}
  \pi_1(\SO(3)) \cong \mathbb{Z}_2.
\end{equation}
For $n\ge 3$, $\pi_1(\SO(n))\cong \mathbb{Z}_2$, and $\Spin(n)$ is the simply connected double cover. A gauge theory with gauge group $\SO(3)$ is not the same global theory as one with gauge group $\SU(2)$, even though the Lie algebras agree.

\subsubsection{\texorpdfstring{Fundamental group of $\U(n)$ and $\SU(n)$}{Fundamental group of U(n) and SU(n)}}

The short exact sequence
\begin{equation}
  1 \to \SU(n) \to \U(n) \xrightarrow{\det} \U(1) \to 1
\end{equation}
yields a fibration $\SU(n) \hookrightarrow \U(n) \to \U(1)$. The associated long exact sequence in homotopy gives
\begin{equation}
  \pi_1(\SU(n)) \to \pi_1(\U(n)) \to \pi_1(\U(1)) \to \pi_0(\SU(n)).
\end{equation}
Since $\pi_0(\SU(n))=0$, $\pi_1(\U(1))\cong \mathbb{Z}$, and $\pi_1(\SU(n))=0$ for $n\ge 2$, one obtains
\begin{equation}
  \pi_1(\U(n))\cong \mathbb{Z},
  \qquad
  \pi_1(\SU(n))=0.
\end{equation}
The integer for $\U(n)$ is the stable winding number detected by the determinant.

\subsection{Holonomy and flat connections}

\subsubsection{Parallel transport}

Let $E \to M$ be a vector bundle with connection $\nabla$. Along a path $\gamma : [0,1] \to M$, a section $s(t)$ of $\gamma^*E$ is \emph{parallel} if
\begin{equation}
  \nabla_{\dot\gamma} s = 0.
\end{equation}
Solving this ODE gives a linear isomorphism
\begin{equation}
  P_\gamma : E_{\gamma(0)} \to E_{\gamma(1)},
\end{equation}
called parallel transport. For concatenated paths, $P_{\gamma_2 * \gamma_1}=P_{\gamma_2}\circ P_{\gamma_1}$. For the reverse path, $P_{\overline\gamma}=P_\gamma^{-1}$.

Fix $x\in M$. The \emph{holonomy group} at $x$ is the subgroup of automorphisms of the fiber $E_x$ generated by parallel transport around based loops:
\begin{equation}
  \Hol_x(\nabla) \subset \operatorname{Aut}(E_x).
\end{equation}

\subsubsection{Flat connections and representations}

A connection is \emph{flat} if $F_\nabla=0$. Flatness implies that locally the connection is gauge equivalent to the trivial connection. Globally, flat connections may still be nontrivial because of holonomy.

If $\nabla$ is flat on a principal $G$-bundle, parallel transport around loops depends only on homotopy class. Hence a flat connection determines a representation
\begin{equation}
  \rho_\nabla : \pi_1(M,x) \to G,
\end{equation}
well-defined up to conjugation. Conversely, a representation of the fundamental group determines a flat connection up to gauge equivalence:
\begin{equation}
  \mathcal{M}_{\mathrm{flat}}(M,G)
  \cong
  \Hom(\pi_1(M),G)/G.
\end{equation}
This identification is fundamental: the classical solutions of Chern--Simons theory are flat connections, so the moduli space of classical solutions is a representation variety.

\subsubsection{Path-ordered exponentials and Wilson loops}

In local coordinates, parallel transport along $\gamma$ is formally
\begin{equation}
  P_\gamma(A)=\mathcal{P}\exp\left(-\int_\gamma A\right),
\end{equation}
where $\mathcal{P}$ denotes path ordering. Under gauge transformation $A \mapsto g^{-1}Ag+g^{-1}\dd g$, holonomy around a loop based at $x$ is conjugated by $g(x)$. Conjugacy-invariant functions of holonomy are gauge invariant. Hence Wilson loops
\begin{equation}
  W_R(\gamma)=\Tr_R P_\gamma(A)
\end{equation}
are natural observables.

\subsection{Symplectic geometry and moduli spaces}

Symplectic geometry enters TQFT because the moduli space of flat connections on a surface is symplectic, and quantization of that space produces the Hilbert space of the theory.

\subsubsection{Symplectic manifolds}

A \emph{symplectic structure} on a manifold $M$ is a closed nondegenerate $2$-form
\begin{equation}
  \omega \in \Omega^2(M),
  \qquad
  \dd\omega=0.
\end{equation}
Nondegeneracy means that at each point the map $T_pM \to T_p^*M$, $v \mapsto \iota_v\omega$ is an isomorphism. In particular, $\dim M$ must be even. The $n$th exterior power satisfies $\omega^n \neq 0$, so a symplectic manifold carries a canonical volume form $\frac{1}{n!}\omega^n$.

\subsubsection{Darboux theorem}

Around every point there are local coordinates $(q_1,p_1,\dots,q_n,p_n)$ such that
\begin{equation}
  \omega = \sum_{j=1}^n \dd q_j \wedge \dd p_j.
\end{equation}
Unlike Riemannian geometry, symplectic geometry has no local invariants: all local structure is the same. The interesting information is global. This is exactly parallel to TQFT: the theory is about global organization, not local metric structure.

\subsubsection{Hamiltonian vector fields and Poisson brackets}

Given $H\in C^\infty(M)$, its Hamiltonian vector field $X_H$ is defined by
\begin{equation}
  \iota_{X_H}\omega = -\dd H.
\end{equation}
The Poisson bracket of functions $F,G$ is
\begin{equation}
  \{F,G\}=\omega(X_F,X_G).
\end{equation}
This gives $C^\infty(M)$ the structure of a Lie algebra, and $X_{\{F,G\}}=[X_F,X_G]$.

\subsubsection{Moment maps and symplectic reduction}

Suppose a Lie group $G$ acts on $(M,\omega)$ preserving $\omega$. For each $\xi\in\mathfrak g$, let $\xi_M$ be the generating vector field. The action is \emph{Hamiltonian} if there exists a map $\Phi : M \to \mathfrak g^*$ such that for every $\xi\in\mathfrak g$,
\begin{equation}
  \dd\langle \Phi,\xi\rangle = -\iota_{\xi_M}\omega.
\end{equation}
The map $\Phi$ is the \emph{moment map}. It packages the conserved quantities associated to the group action.

If $\mu\in\mathfrak g^*$ is a regular value of the moment map, the reduced space
\begin{equation}
  M_{\mu}=\Phi^{-1}(\mu)/G_\mu
\end{equation}
has a natural symplectic form $\omega_\mu$. This is the Marsden--Weinstein--Meyer symplectic reduction theorem.

\subsubsection{The moduli space of flat connections as a symplectic quotient}

Let $\Sigma$ be an oriented closed surface and let $\mathcal{A}$ be the affine space of connections on a fixed principal $G$-bundle over $\Sigma$. The tangent space at any point is $\Omega^1(\Sigma;\mathfrak g)$. The $2$-form
\begin{equation}
  \Omega(\alpha,\beta)=\int_\Sigma \tr(\alpha\wedge\beta)
\end{equation}
defines a symplectic structure on $\mathcal{A}$. The gauge group acts on $\mathcal{A}$, and the moment map is the curvature $\mu(A)=F_A$. Therefore the symplectic quotient is
\begin{equation}
  \mu^{-1}(0)/\mathcal{G}
  =
  \{A : F_A=0\}/\mathcal{G},
\end{equation}
the moduli space of flat connections. This observation, due to Atiyah and Bott, is the central mathematical reason symplectic geometry appears in Chern--Simons theory: the classical phase space on a spatial surface $\Sigma$ is a moduli space of flat connections, and that moduli space is symplectic.

\subsection{Bordism categories and functorial TQFT}

The functorial definition of TQFT, due to Atiyah and Segal, packages the gluing law of path integrals into category theory.

\subsubsection{Bordism category}

Fix a tangential structure, for example orientation. The objects of the bordism category $\Bord_n$ are closed $(n-1)$-manifolds. A morphism $M : \Sigma_0 \to \Sigma_1$ is an $n$-dimensional bordism from $\Sigma_0$ to $\Sigma_1$, meaning an $n$-manifold whose boundary is identified with
\begin{equation}
  \partial M \cong \overline{\Sigma_0} \sqcup \Sigma_1.
\end{equation}
Composition is by gluing along the common boundary. The monoidal structure is disjoint union.

\subsubsection{Definition of TQFT}

An $n$-dimensional TQFT is a symmetric monoidal functor
\begin{equation}
  Z : \Bord_n \to \Vect_k.
\end{equation}
Thus to each closed $(n-1)$-manifold $\Sigma$ one assigns a vector space $Z(\Sigma)$; to each bordism $M : \Sigma_0 \to \Sigma_1$ one assigns a linear map $Z(M): Z(\Sigma_0) \to Z(\Sigma_1)$; disjoint unions go to tensor products; gluing of bordisms goes to composition of linear maps. This formalizes the path-integral slogan that amplitudes compose under gluing.

\subsubsection{Low-dimensional classification}

Oriented $2$-dimensional TQFTs are equivalent to commutative Frobenius algebras. If $Z : \Bord_2 \to \Vect_k$, then the vector space $A=Z(S^1)$ comes equipped with multiplication, unit, comultiplication, and counit obtained from the pair-of-pants bordisms and disks, and these satisfy the Frobenius relations. Conversely, a commutative Frobenius algebra determines a $2$D TQFT. This is a model of how topological gluing laws become algebraic identities.

\subsection{Chern--Simons theory as a TQFT}

\subsubsection{Classical action and flat connections}

On a closed oriented $3$-manifold $M$ with principal $G$-bundle and connection $A$, the Chern--Simons action is
\begin{equation}
  S_{\CS}[A]
  =
  \frac{k}{4\pi}
  \int_M \tr\!\left(A\wedge \dd A + \frac{2}{3}A\wedge A\wedge A\right).
\end{equation}
The Euler--Lagrange equation is $F_A=0$. Classical solutions are flat connections. The classical moduli space is
\begin{equation}
  \mathcal{M}_{\mathrm{cl}}(M,G)
  =
  \{A : F_A=0\}/\mathcal{G}
  \cong
  \Hom(\pi_1(M),G)/G.
\end{equation}
The topology of the manifold enters immediately via $\pi_1(M)$.

\subsubsection{Gauge transformations and level quantization}

For a gauge transformation $g : M \to G$, the Chern--Simons form changes by an exact term plus a global term involving $\tr(g^{-1}\dd g)^{\wedge 3}$. On a closed manifold the exact term integrates to zero, leaving
\begin{equation}
  S_{\CS}[A^g]-S_{\CS}[A]
  =2\pi k\, n(g),
\end{equation}
where $n(g)\in\mathbb{Z}$ is the winding number determined by $[g]\in\pi_3(G)$. The path integral phase $e^{iS_{\CS}}$ is gauge invariant precisely when $k\in\mathbb{Z}$, assuming the standard normalization. The topological input is exactly the identification $\pi_3(G)\cong \mathbb{Z}$.

\subsubsection{Canonical quantization on a surface}

Let spacetime be $\mathbb{R}\times\Sigma$, where $\Sigma$ is a closed oriented surface. The field equation says the spatial connection must be flat, so the reduced classical phase space is the moduli space of flat $G$-connections on $\Sigma$. Quantizing that symplectic manifold produces the Hilbert space assigned by the TQFT to the surface:
\begin{equation}
  \mathcal{H}_\Sigma = Z(\Sigma).
\end{equation}
Topology and symplectic geometry are not parallel stories. They are the classical and quantum halves of the same construction.

\subsubsection{Wilson lines and holonomy}

Given a loop $\gamma$ in a representation $R$, the Wilson operator is
\begin{equation}
  W_R(\gamma)=\Tr_R \mathcal{P}\exp\left(\oint_\gamma A\right).
\end{equation}
Since the classical solutions are flat, Wilson loops are controlled by holonomy classes. In abelian Chern--Simons theory, correlators of Wilson loops reduce to linking numbers; in nonabelian theory they lead to quantum-group and knot-polynomial structures after quantization.

\subsection{Summary}

The conceptual chain is:
\begin{enumerate}
\item A connection on a principal bundle gives parallel transport.
\item For flat connections, parallel transport descends to a representation of $\pi_1$.
\item The moduli space of flat connections on a surface is a symplectic quotient.
\item Quantizing that symplectic quotient gives the state space of Chern--Simons theory.
\item Large gauge transformations are classified by $\pi_3(G)$, forcing level quantization.
\item Wilson lines probe holonomy, and their correlators detect topological linking.
\item Functorially, the entire theory is encoded by a symmetric monoidal functor from bordisms.
\end{enumerate}
This chain is the skeleton of low-dimensional gauge theory.
